\section{Strong Markov Property}

\begin{definition}[Stopping Time]
    A random variable \(T\) with values on \(\NN_0 \cup \Bqty{\infty}\) is a \term{stopping time} with respect to a random process \(\pqty{X_n}_{n \geq 0}\) , if the event \(\Bqty{T = n}\) only depends on \(\pqty{X_0, \ldots, X_n}\), for all \(n \geq 0\).
\end{definition}

\begin{example}[First Hitting Times and Last Exit Time]
    The first hitting times \emph{are} stopping times.

    Recall from definition that \(T_A = \inf\set{n \geq 0}{X_n \in A}\), so
    \[
        \Bqty{T_A = n} = \Bqty{X_0 \notin A, \ldots, X_{n - 1} \notin A, X_n \in A}
    \]
    but the \term{last exit time} \(S_A = \sup \Bqty{n \geq 0}{X_n \in A}\) is \emph{not} a stopping time.
\end{example}

\begin{theorem}[Strong Markov Property]
    Suppose \(X = \pqty{X_n}_{n \geq 0}\) be a Markov chain with \(X \sim \Markov\pqty{\lambda, P}\), and let \(T\) be a stopping time, with \(i \in I\). Then, conditional on the event
    \[
        \Bqty{T < \infty} \cap \Bqty{X_T = i},
    \]
    the process \(\pqty{X_{T + n}}_{n \geq 0} \sim \Markov\pqty{\kronecker_i, P}\), and it is independent of \(\pqty{X_0, X_1. \ldots, X_T}\).
\end{theorem}

\begin{proof}
    We need to show that, for any \(x_0, x_1, \ldots, x_n \in I\) and any finite-length path \(w \in \bigcup_{k \geq 0} I^k\), we have
    \begin{align*}
        &\phantom{=} \Prob\pqty{\Conditional{\pqty{X_0, \ldots, X_T} = w, X_T = x_0, X_{T + 1} = x_1, \ldots, X_{T + n} = x_n}{T < \infty, X_T = i}} \\
        &= \indicator\Bqty{x_0 = i} P\pqty{x_0, x_1} \cdots P\pqty{x_{n - 1}, x_n} \Prob\pqty{\Conditional{\pqty{X_0, \ldots, X_T} = w}{T < \infty, X_T = i}}.
    \end{align*}

    For a fixed \(w \in I^{k + 1}\) and fixed \(x_0, x_1, \ldots, x_n \in I\), we have
    \begin{align*}
        &\phantom{=} \Prob\pqty{\Conditional{\pqty{X_0, \ldots, X_k} = w, X_T = x_0, \ldots, X_{T + n} = x_n}{T < \infty, X_T = i}} \\
        &= \frac{\Prob\pqty{\pqty{X_0, \ldots, X_k} = w, X_T = x_0, \ldots, X_{T + n} = x_n, T = k, X_k = i}}{\Prob\pqty{T < \infty, X_T = i}}.
    \end{align*}
    and for the numerator, we notice that it is equal to
    \begin{align*}
        &\phantom{=} \Prob\pqty{\pqty{X_0, \ldots, X_k} = w, X_k = i, T = k} \cdot \Prob\pqty{\Conditional{X_{k + 1} = x_1, \ldots, X_{k + n} = x_n}{\pqty{X_0, \ldots, X_k} = w, X_k = i, T = k}} \\
        &= \Prob\pqty{\Conditional{X_k = x_0, \ldots, X_{k + n} = x_n}{X_k = i}} \cdot \Prob\pqty{\pqty{X_0, \ldots, X_k} = w, X_k = i, T = k} \\
        &= \indicator \Bqty{x_0 = i} P\pqty{x_0, x_1} \cdots P\pqty{x_{n - 1}, x_n} \Prob\pqty{\Conditional{\pqty{X_0, \ldots, X_T} = w}{T < \infty, X_T = i}}
    \end{align*}
    as desired.
\end{proof}

\begin{example}[Simple Symmetric Random Walk on \(\NN_0\)]
    Recall the setup in \autoref{ex:random-walk-hitting-probability}. We now consider the case where \(p = q\). From previously, we have
    \[
        h_1^A = \frac{1}{2} + \frac{1}{2} h_2^A.
    \]

    Starting at \(2\), let \(\tilde{T}_0\) denote the first time to visit \(0\) under \(\Prob_1\), independent of \(T_1\). Then, by the strong Markov property, conditional on \(\Bqty{T_1 < \infty}\), under \(\Prob_2\), we have \(T_0 = T_1 + \tilde{T}_0\), and hence
    \begin{align*}
        h_2^A &= \Prob_2 \pqty{T_0 < \infty} \\
        &= \Prob_2 \pqty{T_0 < \infty, T_1 < \infty} \\
        &= \Prob_0 \pqty{T_1 < \infty} \Prob\pqty{\Conditional{T_1 + \tilde{T}_0 < \infty}{T_1 < \infty}} \\
        &= \Prob_1 \pqty{T_0 < \infty} \Prob_1 \pqty{T_0 < \infty} \\
        &= h_1^A \cdot h_1^A = \pqty{h_1^A}^2.
    \end{align*}

    So \(h_1^A = \frac{1}{2} + \frac{1}{2} \pqty{h_1^A}^2\), which gives \(h_1^A = h_2^A = 1\).
\end{example}